\documentclass[11pt,a4paper]{article}
\usepackage[margin=1in]{geometry}
\usepackage{amsmath,amssymb,amsthm}
\usepackage{algorithm}
\usepackage{algpseudocode}
\usepackage{booktabs}
\usepackage{hyperref}

\theoremstyle{definition}
\newtheorem{definition}{\textbf{Definition}}[section]
\newtheorem{theorem}{\textbf{Theorem}}[section]
\newtheorem{lemma}[theorem]{\textbf{Lemma}}
\newtheorem{remark}{\textbf{Remark}}[section]

\title{\textbf{Iterative Magnitude Pruning (IMP) and the Lottery Ticket Hypothesis with Weight Rewinding}}
\author{\textbf{Team Scaibu} \\ \small Email: \texttt{jaydeep.9664920749@gmail.com} \\ \small \textbf{Architecture and Core Engine Team}}
\date{\textbf{October 2026}}

\begin{document}
\maketitle

\section{Definitions and Cognitive Mental Models}

\subsection{Formal Mathematical Definition}
\textbf{Definition (Lottery Ticket Hypothesis and Iterative Magnitude Pruning):}
Modern deep neural networks contain billions of parameters, exhibiting massive overparameterization. The \textbf{Lottery Ticket Hypothesis (LTH)} (Frankle \& Carbin, 2019) conjectures that a randomly initialized, dense feedforward network $f(x; \theta_0)$ with $\theta_0 \sim \mathcal{D}_\theta$ contains a sparse subnetwork $f(x; m \odot \theta_0)$, parameterized by a binary mask $m \in \{0, 1\}^P$ ($m \odot \theta \triangleq (m_1 \theta_1, \dots, m_P \theta_P)$), that can be trained in isolation to achieve test accuracy comparable to the original dense network within at most the same number of training iterations.

The sparse mask $m$ is discovered via \textbf{Iterative Magnitude Pruning (IMP)} with weight rewinding:
\begin{enumerate}
    \item Train the dense network from initialization $\theta_0$ (or early warmup checkpoint $\theta_k$, $k \ge 0$) to convergence $\theta_T$.
    \item Compute parameter saliencies by absolute magnitude: $s_i \triangleq |\theta_{T, i}|$.
    \item Determine the pruning cutoff threshold $\tau_s$ at desired sparsity quantile $s \in [0, 1)$:
    {\boldmath
    \begin{equation}
    \tau_s \triangleq \inf\left\{ \tau \in \mathbb{R}_+ : \frac{1}{P} \sum_{i=1}^P \mathbf{1}\left( |\theta_{T, i}| \le \tau \right) \ge s \right\}
    \end{equation}
    }
    \item Construct the binary survival mask $m \in \{0, 1\}^P$:
    {\boldmath
    \begin{equation}
    m_i \triangleq \mathbf{1}\left( |\theta_{T, i}| > \tau_s \right)
    \end{equation}
    }
    \item \textbf{Rewind}: Reset the surviving parameters to their original values at step $k$: $\theta_{\text{pruned}} = m \odot \theta_k$.
\end{enumerate}

\subsection{Expanded Non-Technical Definition (Intuition for Anyone)}
\begin{enumerate}
    \item \textbf{Buying 1,000 Lottery Tickets}: When you initialize a massive AI model with 100 billion parameters, you are buying millions of lottery tickets. Most connections are useless, but a few lucky combinations of initialized weights happen to form an extraordinarily capable sparse circuit.
    \item \textbf{Pruning the Losers}: By letting the full network train, the "winning tickets" grow large and strong, while the useless tickets remain tiny. Pruning chops off the smallest 80\% of weights, leaving only the winning skeleton.
    \item \textbf{The Magic of Rewinding}: If you take that sparse skeleton and re-initialize it with random new numbers, it fails to train. But if you rewind it back to the exact initial numbers it had on Day 1 ($\theta_0$), it trains to full performance!
\end{enumerate}

\subsection{Child-Friendly Mental Model (Physical Metaphor)}
Imagine carving a marble sculpture:
\begin{itemize}
    \item \textbf{Dense Initialization}: A huge, heavy block of raw stone.
    \item \textbf{Training and Pruning}: The sculptor chisels away 80\% of the unnecessary stone, revealing the hidden statue inside.
    \item \textbf{Lottery Ticket}: The statue was always inside the original block; the training process simply showed the sculptor where to cut!
\end{itemize}

\section{Core Mathematical Equations and Definitions}

\subsection{Masked Loss Landscape Optimization}
{\boldmath
\begin{equation}
\min_{m \in \{0, 1\}^P} \mathcal{L}\left( m \odot \theta^* \right) \quad \text{subject to} \quad \frac{\|m\|_0}{P} \le 1 - s
\end{equation}
}
\begin{itemize}
    \item \textbf{Formal Definition}: Combinatorial sparsity-constrained empirical risk minimization over binary hypercube $\{0, 1\}^P$.
    \item \textbf{What It Means}: Finds the smallest subnetwork that preserves model training dynamics.
    \item \textbf{Why It Is Important}: Connects structural parameter redundancy to downstream inference compression.
\end{itemize}

\subsection{Weight Rewinding Parameter Equation}
{\boldmath
\begin{equation}
\theta_{\text{sub}}^{(t=0)} = m \odot \theta_k, \quad \text{where} \quad m_i = \begin{cases} 1 & \text{if } |\theta_{T, i}| \ge \tau_s \\ 0 & \text{if } |\theta_{T, i}| < \tau_s \end{cases}
\end{equation}
}
\begin{itemize}
    \item \textbf{Formal Definition}: Initialization state of the pruned winning ticket subnetwork rewound to step $k$.
    \item \textbf{What It Means}: Preserves the precise random seed alignment that enabled SGD convergence while setting pruned weights strictly to zero.
    \item \textbf{Why It Is Important}: In large LLMs, rewinding to step $k=0$ often fails due to chaos/instability; rewinding to step $k \approx 0.05 T$ (linear stability checkpoint) restores full winning ticket recoverability.
\end{itemize}

\section{Rigorous Mathematical Analysis Checklist}

\begin{itemize}
    \item \textbf{Notation}: $P$ (total parameters), $\theta_0$ (initial weights), $\theta_T$ (trained weights), $m \in \{0, 1\}^P$ (binary mask), $s \in [0, 1)$ (target sparsity).
    \item \textbf{Epistemic Status}: Empirically established in deep learning theory (Frankle \& Carbin, 2019; Malach et al., 2020).
    \item \textbf{Dependency Graph}: Trained Weights $\to$ Absolute Magnitude Ranking $\to$ Threshold Cutoff $\to$ Binary Mask $\to$ Weight Rewinding.
    \item \textbf{Dimensions}: Weight vectors $\theta \in \mathbb{R}^P$; mask $m \in \{0, 1\}^P$; scalar sparsity $s \in [0, 1)$.
    \item \textbf{Regimes}: Unstructured sparsity across convolutional networks, vision transformers, and large language models.
    \item \textbf{Invariants}: Exact conservation of surviving parameter count: $\|m\|_0 = \lceil (1 - s) P \rceil$.
    \item \textbf{Isomorphisms}: Isomorphic to subset selection in LASSO regression and hard thresholding pursuit.
\end{itemize}

\section{Theoretical Theorems and Formal Proofs}

\begin{theorem}[Exact Sparsity Guarantee of Magnitude Quantile Pruning]
Let $\theta \in \mathbb{R}^P$ be a continuous parameter vector with distinct absolute values $|\theta_1| < \dots < |\theta_P|$. For any target sparsity $s \in [0, 1)$, let $K = \lfloor s \cdot P \rfloor$ denote the number of weights to prune. If the pruning threshold is chosen as $\tau = |\theta_{\pi(K)}|$ where $\pi$ is the sorting permutation $|\theta_{\pi(1)}| \le \dots \le |\theta_{\pi(P)}|$, the achieved empirical sparsity $s_{\text{achieved}} \triangleq \frac{1}{P} \sum_{i=1}^P (1 - m_i)$ satisfies:
{\boldmath
\begin{equation}
s_{\text{achieved}} = \frac{\lfloor s \cdot P \rfloor}{P} \in \left[ s - \frac{1}{P}, \, s \right]
\end{equation}
}
\end{theorem}

\begin{proof}
By sorting indices by absolute values, the set of pruned indices is defined as:
\begin{equation}
\mathcal{I}_{\text{prune}} = \{ \pi(1), \pi(2), \dots, \pi(K) \}
\end{equation}
where $K = \lfloor s \cdot P \rfloor$.
By definition of the binary mask:
\begin{equation}
m_i = \begin{cases} 0 & \text{if } i \in \mathcal{I}_{\text{prune}} \\ 1 & \text{if } i \notin \mathcal{I}_{\text{prune}} \end{cases}
\end{equation}
The number of zeroed elements is $|\mathcal{I}_{\text{prune}}| = K$.
The achieved sparsity is:
\begin{equation}
s_{\text{achieved}} = \frac{1}{P} \sum_{i=1}^P (1 - m_i) = \frac{|\mathcal{I}_{\text{prune}}|}{P} = \frac{\lfloor s \cdot P \rfloor}{P}
\end{equation}
Since $s \cdot P - 1 < \lfloor s \cdot P \rfloor \le s \cdot P$, dividing by $P$ yields:
\begin{equation}
s - \frac{1}{P} < s_{\text{achieved}} \le s
\end{equation}
As $P \to \infty$, $s_{\text{achieved}} \to s$ exactly, ensuring the desired sparsity constraint is strictly honored.
\end{proof}

\begin{theorem}[First-Order Loss Taylor Perturbation from Magnitude Pruning]
Let $\mathcal{L}(\theta)$ be a twice continuously differentiable loss function. Expanding $\mathcal{L}$ about the trained parameter point $\theta_T$ where $\nabla \mathcal{L}(\theta_T) \approx \mathbf{0}$, the increase in loss resulting from zeroing pruned parameters $\Delta \theta = m \odot \theta_T - \theta_T$ is bounded by:
{\boldmath
\begin{equation}
\Delta \mathcal{L} \approx \frac{1}{2} \sum_{i \in \mathcal{I}_{\text{prune}}} H_{ii} \theta_{T, i}^2 + \mathcal{O}\left( \max_{i \in \mathcal{I}_{\text{prune}}} |\theta_{T, i}|^3 \right)
\end{equation}
}
where $H = \nabla^2 \mathcal{L}(\theta_T)$ is the Hessian matrix.
\end{theorem}

\begin{proof}
Consider the multivariate Taylor expansion of $\mathcal{L}(\theta_T + \Delta \theta)$ about $\theta_T$:
\begin{equation}
\mathcal{L}(\theta_T + \Delta \theta) = \mathcal{L}(\theta_T) + \nabla \mathcal{L}(\theta_T)^\top \Delta \theta + \frac{1}{2} \Delta \theta^\top H \Delta \theta + \mathcal{O}(\|\Delta \theta\|_2^3)
\end{equation}
At local convergence, $\nabla \mathcal{L}(\theta_T) \approx \mathbf{0}$.
For the pruning perturbation, $\Delta \theta_i = -\theta_{T, i}$ for $i \in \mathcal{I}_{\text{prune}}$ and $\Delta \theta_i = 0$ for $i \notin \mathcal{I}_{\text{prune}}$.
Substituting into the quadratic form:
\begin{equation}
\Delta \theta^\top H \Delta \theta = \sum_{i \in \mathcal{I}_{\text{prune}}} \sum_{j \in \mathcal{I}_{\text{prune}}} H_{ij} \theta_{T, i} \theta_{T, j}
\end{equation}
Under the standard diagonal dominance approximation of the Hessian ($H_{ij} \ll H_{ii}$ for $i \neq j$):
\begin{equation}
\Delta \mathcal{L} \approx \frac{1}{2} \sum_{i \in \mathcal{I}_{\text{prune}}} H_{ii} \theta_{T, i}^2
\end{equation}
Because magnitude pruning explicitly selects the smallest values of $|\theta_{T, i}| \le \tau_s$, it minimizes the dominant scalar multiplier $\theta_{T, i}^2$, thereby minimizing the first-order destruction of model loss compared to random or sign-based pruning.
\end{proof}

\section{Architectural Questions and Real-World Trade-Offs}

\subsection{Architectural Question 1: Unstructured Sparsity on GPU Hardware}
\textbf{Question:} Why doesn't 50\% unstructured magnitude pruning yield a $2\times$ inference speedup on standard GPUs?\\
\textbf{Answer:} Standard GPU tensor cores (e.g., NVIDIA H100) execute dense Matrix-Matrix Multiplications (GEMM) using aligned systolic memory layouts. Randomly setting scalar weights to zero creates irregular memory access patterns and branch divergence. Without specialized sparse accelerators or structured patterns (like 2:4 sparsity or block sparsity), unstructured sparse matrices actually run slower than dense GEMM due to memory layout overhead.

\subsection{Architectural Question 2: Why Rewind Instead of Fine-Tuning?}
\textbf{Question:} What is the practical difference between fine-tuning pruned weights versus rewinding them?\\
\textbf{Answer:} In standard compression (prune-and-finetune), weights are kept at $\theta_T$ and fine-tuned for a few epochs. This works well for modest sparsity ($<50\%$). However, at extreme sparsity ($>80\%$), fine-tuning gets stuck in damaged local minima. Rewinding to early training checkpoint $\theta_k$ allows the sparse subnetwork to re-traverse the full loss optimization trajectory, discovering alternative basin solutions that match dense accuracy.

\section{Algorithmic Implementation Specification}

\begin{algorithm}[H]
\caption{Iterative Magnitude Pruning with Weight Rewinding}
\begin{algorithmic}[1]
\Require Current weights $\boldsymbol{\theta} \in \mathbb{R}^P$, initial weights $\boldsymbol{\theta}_0 \in \mathbb{R}^P$, target sparsity $s \in [0, 1)$, rewind flag $\mathcal{R}$
\Ensure Pruned weights $\boldsymbol{\theta}' \in \mathbb{R}^P$, binary mask $\mathbf{m} \in \{0, 1\}^P$, achieved sparsity $s_{\text{ach}}$, threshold $\tau$, surviving count $P_{\text{surv}}$
\State $P \gets \operatorname{len}(\boldsymbol{\theta})$
\If{$s = 0$}
    \State \Return $(\boldsymbol{\theta}_0 \text{ if } \mathcal{R} \text{ else } \boldsymbol{\theta}, [1]^P, 0.0, 0.0, P)$
\EndIf
\State $K \gets \operatorname{clip}(\lfloor s \cdot P \rfloor, 0, P - 1)$
\State Sort parameter indices by magnitude: $\boldsymbol{\pi} \gets \operatorname{argsort}([|\theta_i| \text{ for } i \in [0, P-1]])$
\State $\mathcal{I}_{\text{prune}} \gets \operatorname{Set}(\boldsymbol{\pi}[0 : K])$
\State $\tau \gets |\boldsymbol{\theta}[\boldsymbol{\pi}[K]]|$
\State Initialize $\mathbf{m} \gets [\,], \quad \boldsymbol{\theta}' \gets [\,], \quad P_{\text{surv}} \gets 0$
\For{$i = 0$ \textbf{to} $P - 1$}
    \If{$i \in \mathcal{I}_{\text{prune}}$}
        \State Append $0$ to $\mathbf{m}$
        \State Append $0.0$ to $\boldsymbol{\theta}'$
    \Else
        \State Append $1$ to $\mathbf{m}$
        \State Append $\boldsymbol{\theta}_0[i]$ if $\mathcal{R}$ else $\boldsymbol{\theta}[i]$ to $\boldsymbol{\theta}'$
        \State $P_{\text{surv}} \gets P_{\text{surv}} + 1$
    \EndIf
\EndFor
\State $s_{\text{ach}} \gets (P - P_{\text{surv}}) / P$
\State \Return $(\boldsymbol{\theta}', \mathbf{m}, s_{\text{ach}}, \tau, P_{\text{surv}})$
\end{algorithmic}
\end{algorithm}

\section{Big-$\Theta$ Complexity Analysis}

\begin{itemize}
    \item \textbf{Magnitude Extraction Pass}: $\Theta(P)$ absolute value evaluations.
    \item \textbf{Index Magnitude Sorting}: $\Theta(P \ln P)$ comparison sort pass.
    \item \textbf{Mask Construction and Rewinding}: $\Theta(P)$ sequential array assignment.
    \item \textbf{Total Time Complexity}: $\Theta(P \ln P)$ bounded by the sorting step.
    \item \textbf{Space Complexity}: $\Theta(P)$ auxiliary memory for index tracking and binary mask.
\end{itemize}

\section{Single Canonical Repository Reference}

The complete deterministic scalar implementation, verification suite, and unit test benchmarks are published at:
\begin{center}
\url{https://github.com/Chief-Strategist-J/llm-obs-infra}
\end{center}

\end{document}
